\chapter{MLS Level~3 datasets}
\label{chap:Level3}

\section{Introduction}

The other chapters of this guide describe MLS ``Level~2'' data --
geophysical products along the instrument observing track.  This
chapter describes the MLS ``Level~3'' -- geophysical products mapped
onto regular grids.  There is a wide array of techniques for
transforming along-track observations into such regular grids,
including ``kriging'' approaches \citep[e.g.,][]{Cressie:2015},
Fourier-series-based approaches \citep{Salby:1982}, approaches using
various windowing functions such as Hamming windows, and simple
binning approaches.  All methods have strengths and weaknesses, and
the choice of which one to use is best made by the individual data
user, given the nature of the specific scientific question to be
tackled.  For this reason, until 2019, the MLS team had not produced
any ``official'' Level~3 data products.

However, in order to facilitate more rapid ``quick-look'' studies by
MLS data users, the MLS team is now producing and distributing a
select set of Level~3 products derived using a simple ``binning''
approach (i.e., simply reporting the mean measured value within a
given spatial/temporal ``box'').  As part of the generation of these
data products, all the data screening rules described in
chapters~\ref{chap:Essential} and~\ref{chap:Results} are applied,
presenting a significant additional simplification to data users.

The Level~3 products can be divided into separate families
\begin{description}
\item[Daily and monthly 3D grids:] Latitude/Longitude/Pressure and
  Latitude/Longitude/Potential-Temperature grids of most products.
  Daily grids are stored in individual files for each product.
  Monthly grids are stored in yearly files containing separate
  entries for each of the 12 months.
\item[Daily and monthly latitude/pressure and
  latitude/potential-temperature zonal means:] Stored in annual files
  for each product, with 365 or 366 entries per year for the daily
  files, and 12 entries per year for the monthly files.
\item[Daily and monthly equivalent latitude/potential temperature
  means:] Stored as for the zonal mean products above.
\item[Daily and monthly vortex average timeseries:] Annual files
  containing daily and monthly profiles of vortex averages of selected
  species on a potential temperature vertical coordinate.
\end{description}
For species having a non-negligible diurnal variation in all or part
of their vertical range, the files contain separate averages of
\begin{itemize}[itemsep=0.5\baselineskip]
\item all the profiles,
\item all profiles measured on the ascending side of the orbit (mainly
  daytime observations),
\item all profiles measured on the descending side of the orbit (mainly
  nighttime observations),
\item all observations made with solar zenith angle less than
  90\textdegree{} (i.e., daytime observations).
\item all observations made with solar zenith angle greater than
  110\textdegree{} (i.e., nighttime observations).
\end{itemize}

The MERRA-2 reanalysis temperature fields (interpolated to each MLS
observation point) are used to convert from pressure to potential
temperature surfaces.  Additionally MERRA-2 fields are used to
construct equivalent latitude and vortex average information.  The
interpolated MERRA-2 fields used are taken from the MLS ``Derived
Meteorological Products'' \citep{ManneyEtal:2007}.  Note that, for
quantities on potential temperature surfaces, the MLS Level~2 profiles
on pressure levels are first interpolated (from $\log(p)$ to
$\log(\theta)$) to the fixed potential temperature surfaces before
being binned (as opposed to directly binning Level~2 data points in
potential temperature windows).  These interpolated values (and the
interpolated \texttt{L2GPPrecision} field are then used to construct
the \texttt{minimum}, \texttt{maximum}, \texttt{std\_dev}, and
\texttt{n\_values} fields. \mycomment{Ryan, is that correct,
  particularly the n\_values part?}

\section{Level~3 data files}

All of the MLS Level 3 data files described here are available from
the NASA Goddard Space Flight Center Data and Information Services
Center (GSFC-DISC, \url{http://disc.gsfc.nasa.gov/}).  These data
products are stored in standard CF-compliant NetCDF-4 files.  With the
ease of reading NetCDF-4 files, a dedicated reader has not been made
available at this time.

Two different temporal resolutions of Level~3 data are provided—daily
(\texttt{L3SD}), and monthly (\texttt{L3SM}). These files contain all
six types of fields: latitude/longitude grids on pressure and
potential temperature, zonal means on both pressure and potential
temperature vertical coordinates, means around equivalent latitude
``circles'' on potential temperature, and polar vortex averages as
function of potential temperature.  The \texttt{L3SD} files are one
file per day per Level~2 data product and the \texttt{L3SM} files are
one file per year per Level~2 data product (12 time entries in the
file).  To aid in the creation of long-term datasets, there are also
the \texttt{L3SY} files which rollup all of the \texttt{L3SD} data
fields except for the daily gridded (latitude/longitude files), which
are not included for file size considerations.  The \texttt{L2SY}
files contain(365 or 366 time entries as appropriate.

The filenames for these products using the following conventions which
are similar varying in the date format:
\begin{quote}
\verb+MLS-Aura_L3SD-<product>_v04-20-c01_<yyyy>d<ddd>.nc+\\
\verb+MLS-Aura_L3<type>-<product>_v04-20-c01_<yyyy>.nc+
\end{quote}
Here, \verb+<type>+ is \texttt{SM} or \texttt{SY}, \verb+<product>+
is \texttt{BrO}, \texttt{O3}, \texttt{Temperature}, etc.  The observed
date is the four digit year \verb+<yyyy>+ with the optional three
digit day of year \verb+<ddd>+ (\texttt{001} = 1 January).

Within each file, the data are stored in NetCDF-4 groups.  The daily
and monthly files contain grids, while all files contain zonal mean
and vortex average groups.  These are named as the following where
product is the same as listed in the filename:

\vspace{\baselineskip}\noindent\begin{tabular}{>{\ttfamily}r>{\raggedright}p{4.4in}}
  <product> PressureGrid: & 4$\times$5\textdegree{} geodetic
  longitude/latitude grid on pressure surfaces (\texttt{L3SD} and
  \texttt{L3SM} only)\tabularnewline
  
<product> ThetaGrid: & 4$\times$5\textdegree{} geodetic
  longitude/latitude grid on potential temperature surfaces
  (\texttt{L3SD} and \texttt{L3SM} only)\tabularnewline
  
  <product> PressureZM: & 4\textdegree{} geodetic latitude zonal mean
on pressure surfaces\tabularnewline

<product> ThetaZM: & 4\textdegree{} geodetic latitude zonal mean on
potential temperature surfaces\tabularnewline

<product> EqLZM: & 4\textdegree{} equivalent latitude zonal mean on
potential temperature surfaces\tabularnewline

<product> VortexAvg: & Polar vortex average values for each
hemisphere on potential temperature surfaces.\tabularnewline
\end{tabular}

\vspace{\baselineskip}Diurnal product files have all of these groups
repeated for day, night, ascending profiles only, and descending
profiles only (add \texttt{Day}, \texttt{Night}, \texttt{Asc},
\texttt{Desc}, respectively to the group name). Each group contains
the following data fields:

\vspace{\baselineskip}\noindent\begin{tabular}{>{\ttfamily}r>{\raggedright}p{4.9in}}
  lon: & The longitude at center of the grid cell
  (\texttt{PressureGrid} and \texttt{ThetaGrid} groups
  only)\tabularnewline

  lon\_bnds: & The longitude at the boundaries of the grid cell
  (\texttt{PressureGrid} and \texttt{ThetaGrid} groups
  only)\tabularnewline

  lat: & The latitude at center of the grid cell\tabularnewline

  lat\_bnds: & The latitude at the boundaries of the grid
  cell\tabularnewline

  time: & The date of the grid cell (days since
  1950-01-01)\tabularnewline

  time\_bnds: & The boundaries of the time period (days since
  1950-01-01)\tabularnewline

  lev: & The vertical coordinate (i.e. pressure/potential temperature
  values)\tabularnewline

  value: & The average value for each bin\tabularnewline

  rms\_uncertainty: & The root mean square of all the Level~2
  \texttt{L2GPPrecision} values that contributed to each bin (note
  that this has \emph{not} been divided by
  $\sqrt{\texttt{nvalues}}$\tabularnewline

  minimum: & The minimum value in each bin\tabularnewline

  maximum: & The maximum value in each bin\tabularnewline

  std\_dev: & The standard deviation of the data in each
  bin\tabularnewline

  nvalues: & The number of valid data points found in each
  bin\tabularnewline
\end{tabular}

\section{Guidance for users of Level~3 data}

As discussed above, these Level~3 products are generated using only
the Level~2 data points that meet all the screening criteria described
in the earlier chapters of this document, presenting a significant
simplification for users. In cases where there are no good Level~2
data available to populate a Level~3 data point the \texttt{nvalues}
field is set to zero and the other data fields are set to the
\texttt{\_fillValue} in the NetCDF file.

When considering Level~3 data, as with every geophysical measurement,
attention should be paid to the associated uncertainties.  The
averaging inherent in generating Level~3 data will act to reduce
``precision''-related errors, those due to random noise in the MLS
radiance measurements, by $1/\sqrt{n}$, where $n$ is the number of
Level~2 data points in the average (which can be deduced by
considering the \texttt{nvalues} field).  In contrast, the
``accuracy''-related errors, those due to biases in the MLS
measurement system, will in most cases be propagated through the
averaging process unchanged, and should be assumed to apply to the
Level~3 values in the same manner as they do in Level~2 datasets.

%%% Local Variables:
%%% mode: latex
%%% TeX-master: "v4-x-report"
%%% End:
